\documentclass[a4paper, 12pt, UTF8]{ctexart}

% ============ 宏包 ============
\usepackage[top=2.5cm, bottom=2.5cm, left=2.5cm, right=2.5cm]{geometry}
\usepackage{amsmath, amssymb}
\usepackage{enumitem}
\usepackage{setspace}
\usepackage{hyperref}
\usepackage{indentfirst}

% ============ 格式设置 ============
\linespread{1.3}
\setlength{\parindent}{2em}
\hypersetup{colorlinks=true, linkcolor=blue, citecolor=blue, urlcolor=blue}

% ============ 标题信息 ============
\title{\textbf{JUNO实验反应堆中微子振荡分析与高性能拟合算法研究} \\ \vspace{0.3em} \Large 研究计划}
\author{薛景秦}
\date{}

% ============ 正文 ============
\begin{document}
\maketitle
\thispagestyle{empty}

%% ================================================================
\section{选题背景}

中微子振荡现象——即不同味本征态之间的量子力学转化——是超越标准模型新物理的确凿证据。在三代中微子框架下，振荡行为由Pontecorvo-Maki-Nakagawa-Sakata（PMNS）混合矩阵中的三个混合角（$\theta_{12}$、$\theta_{13}$、$\theta_{23}$）、两个独立质量平方差（$\Delta m^2_{21}$、$\Delta m^2_{31}$）及CP破坏相$\delta$共六个参数描述。其中，太阳振荡参数$\sin^2\theta_{12}$和$\Delta m^2_{21}$的高精度测量尤为关键：它们是未来无中微子双贝塔衰变（$0\nu\beta\beta$）实验确定中微子有效马约拉纳质量及判定中微子本质属性（Majorana或Dirac）的前提条件。因此，将这些参数的测量精度推进至亚百分比量级，是粒子物理与宇宙学交叉领域中最迫切的精密测量任务之一。

江门中微子实验（JUNO）依托地下700\,米处的20\,千吨液体闪烁体探测器，通过逆贝塔衰变（IBD）反应探测52.5\,公里基线处的反应堆反中微子，旨在实现振荡参数的亚百分比量级精密测量。为实现这一目标，JUNO在运行早期需要应对一系列方法学挑战：（1）\textbf{低统计量}将导致重建能谱部分能量箱中事例稀少，泊松涨落的非高斯性质可能引入系统性的参数估计偏差，如何选择合适的能谱分箱方案以有效抑制该偏差是亟待解决的问题；（2）台山反中微子观测站（TAO）作为近端探测器尚未积累充足数据，振荡分析将不得不\textbf{依赖反应堆能谱理论模型}（如Huber-Mueller转换模型或核数据库求和模型），不同模型预言的能谱形状差异——特别是5--7\,MeV区间的精细结构——将直接传导为振荡参数的系统误差，因此需要需要量化反应堆模型依赖引入的系统误差以及制定合理的反应堆中微子能谱的分bin策略；（3）在$\Delta m^2_{31}$的测量中，振荡相位简并可能导致似然面出现多个分离的局域极小值，加之低统计量下似然面的非抛物线结构，\textbf{Wilks定理的正则性条件可能遭到违反}，标准$\Delta\chi^2$阈值将给出覆盖不足的置信区间；（4）为保证统计推断的稳健性而拟采用的Feldman-Cousins（FC）方法需要在参数空间中进行大规模伪实验拟合，\textbf{计算量极为巨大}，亟需开发高性能计算框架以使之具备可行性。

本研究拟围绕上述挑战体系，从数值算法优化、能谱分箱与分段策略、以及严格统计方法三个维度展开系统研究，旨在为JUNO早期运行阶段的精密测量提供完整、可靠的分析方法支持。

%% ================================================================
\section{研究内容}

\subsection{拟开发基于张量化的高性能前向折叠拟合框架}

JUNO反应堆中微子振荡分析采用前向折叠方法：从反应堆端的反中微子发射谱出发，经振荡传播、IBD反应截面加权、沉积能转换、液体闪烁体非线性响应和探测器能量分辨率模糊，最终生成理论预期的重建能谱与实验观测进行比较。该过程在数值实现上可表示为矩阵链运算$\vec{\phi}_{\rm rec} = \mathbf{R} \cdot \mathbf{M} \cdot \vec{\phi}_{\bar\nu}$，其中$\mathbf{M}$为IBD截面映射矩阵，$\mathbf{R}$为探测器响应矩阵。初步分析表明，前向折叠能谱生成过程中，探测器响应矩阵的构造（涉及大量误差函数$\mathrm{erf}$计算）和IBD截面映射（涉及大维度矩阵运算）占据了绝大部分计算时间，而极小化算法本身的开销相对较小。因此，本项目拟将优化重点集中于能谱生成链的加速。

具体而言，拟实施以下三层递进优化策略：（1）\textbf{静态量预计算与动态结果缓存}——将不随拟合参数变化的静态量（裂变同位素基准能谱、IBD截面映射矩阵、非线性曲线、本底模板等）一次性预计算并常驻内存，同时对依赖拟合参数的动态计算结果引入函数级缓存机制，消除重复计算；（2）\textbf{IBD截面映射的稀疏矩阵表示}——利用IBD运动学约束赋予截面映射矩阵的带状稀疏结构，将其从稠密格式转换为压缩稀疏行（CSR）格式存储，从而将矩阵-向量乘法的计算复杂度从$\mathcal{O}(N^2)$降至$\mathcal{O}(Nw)$（$w$为带宽），同时改善缓存局部性；（3）\textbf{响应矩阵的分块构造}——采用列方向分块策略处理大维度探测器响应矩阵，通过改善数据局部性进一步提升$\mathrm{erf}$密集计算的效率。

本项目预期通过上述优化策略，显著缩减单次能谱拟合的计算时间，使Feldman-Cousins方法所需的百万量级拟合在有限计算资源下具备可行性，为后续大规模参数扫描与统计检验研究奠定计算基础。

\subsection{拟研究能谱分箱与分段策略}

\textbf{重建能谱分箱策略}：在JUNO早期低统计量阶段，过细的重建能谱箱宽将导致部分能量箱中事例数极少，泊松分布的非高斯性质变得显著，标准的Neyman或Pearson加权$\chi^2$方案可能引入系统性的参数估计偏差。本项目拟采用Combine-Neyman-Pearson（CNP）组合加权方案以抑制一阶偏差，并拟通过生成大规模伪实验，系统扫描不同箱宽对$\sin^2\theta_{12}$和$\Delta m^2_{21}$测量偏差的影响，确定在低统计量条件下兼顾偏差抑制与物理灵敏度保留的最优箱宽配置。

\textbf{反应堆能谱分段策略}（JUNO与大亚湾联合分析）：在TAO数据充足之前，拟借助大亚湾实验积累的高统计量近端数据对反应堆能谱形状进行约束，将反应堆反中微子能谱在能量轴上分为若干段，各段内的谱形固定为基准模型预言，但段间的相对归一化作为自由参数在拟合中浮动，由大亚湾数据提供约束。分段宽度的选择涉及根本性权衡——过细将导致自由参数过多、大亚湾无法有效约束，可能引发过拟合；过粗则无法容纳局域的模型偏差，导致模型依赖性残留。本项目拟通过构建包含模型不确定性的伪实验研究，系统评估不同分段宽度对振荡参数测量精度与模型依赖引入的系统误差的影响，确定最优的分段方案，以在模型依赖性抑制与物理灵敏度保留之间取得平衡。

\subsection{拟采用Feldman-Cousins方法开展统计推断稳健性研究}

在JUNO早期对$\Delta m^2_{31}$的测量中，振荡概率的周期性可能导致$\chi^2$剖面在参数空间中出现多个分离的局域极小值（相位简并），叠加低统计量下似然面的非抛物线结构，Wilks定理的正则性条件——包括似然函数的单模态性、参数不在边界上、以及充分大的样本量——均可能遭到违反。在此条件下，直接套用标准$\Delta\chi^2$阈值构建置信区间将面临覆盖不足的风险。

为确保统计推断的稳健性，本项目拟采用Feldman-Cousins（FC）方法开展大规模频率学覆盖性检验。具体而言，拟在参数空间的每个假设点上生成大量伪实验数据集，对每个伪实验进行完整的剖面似然比拟合，收集检验统计量$\Delta\chi^2$的经验分布，从中提取参数依赖的临界值以替代固定的Wilks阈值。本项目拟分别对太阳振荡参数（$\Delta m^2_{21}$, $\sin^2\theta_{12}$）开展二维FC检验、对大气振荡参数$\Delta m^2_{31}$开展一维FC检验，系统评估Wilks近似在不同参数和不同曝光量下的适用性，量化可能存在的覆盖偏差，并给出基于FC校正的可靠置信区间。该研究的实施依赖于上述高性能拟合框架提供的计算支撑。

%% ================================================================
\section{预期目标}

\textbf{（1）物理精度目标}：预期通过系统的分箱策略研究，确定JUNO早期低统计量阶段的最优重建能谱箱宽，有效抑制非高斯统计涨落引入的参数估计偏差；通过联合分析中反应堆能谱分段方案的优化，量化并控制反应堆模型依赖性引入的系统误差至可接受水平，确保系统误差不成为精密测量的限制因素。

\textbf{（2）算法性能目标}：预期建立一套兼具高计算效率与高数值精度的通用前向折叠拟合平台，通过张量化运算链、稀疏矩阵表示和分块矩阵构造等优化策略，显著缩减单次能谱拟合的计算时间，使基于Feldman-Cousins方法的大规模伪实验分析在标准计算资源下成为可能。该框架的设计理念预期具有方法学的一般性，可推广至其他采用前向折叠方法的中微子实验谱分析。

\textbf{（3）科学贡献目标}：预期为JUNO首篇物理分析提供一套稳健的基准分析工具链，涵盖高性能拟合器、经优化的能谱分箱与分段策略、以及基于FC方法的严格统计推断方案。通过大规模频率学覆盖性检验，预期系统评估Wilks近似在JUNO各振荡参数测量中的适用性，为低统计量阶段的置信区间构建提供可靠的统计保障，并为合作组制定数据发表策略提供量化的方法学建议。

\end{document}
